\begin{abstract}
随着\"双碳\"战略深入推进，绿电直连型电氢氨一体化园区成为解决新能源消纳和化工行业深度脱碳的重要路径。本文针对绿电直连型电氢氨园区的优化运行问题，建立了从功率平衡分析到储能配置优化的完整数学模型体系。

针对问题一，基于典型日风光出力曲线和常规负荷曲线，建立了逐时功率平衡模型，计算了园区典型日的用电量、发电量、购售电等关键指标。结果表明：在36 t/d产能下，自发自用比例(28.16\%)和上网电量比例(35.92\%)未达到绿电直连要求，主要原因是制氢氨负荷(20.75 MW)远小于新能源装机容量(104 MW)，大量绿电上网导致自用比例不足。

针对问题二，建立了基于边际成本排序的离散调度优化模型，在72 t/d产能下，对各产量水平(72-36 t/d)分别优化了最优生产时段安排。结果表明：72 t/d、63 t/d和54 t/d可满足全部绿电指标，45 t/d及以下部分指标不达标。24种风光场景下的统计分析显示，72 t/d时18/24场景全满足，而36 t/d时0/24场景全满足。

针对问题三，建立了连续功率可调(10\%-100\%)的线性规划模型。相比离散调节，连续调节显著改善了绿电指标：各产量下自用比例保持在82\%以上，上网比例控制在9\%以内，全达标场景数从Q2的0-18个提升至Q3的18-22个，连续调节在低产量下优势尤为显著。

针对问题四，分析了离网运行模式下园区的能源自给状况，24种场景下平均产能利用率37.7\%，平均可再生利用率93.7\%。储能配置优化的Pareto分析表明，最优储能容量为20 MWh，可略微提升产量约0.2\%。相比联网模式，离网模式吨氨成本较低但产能大幅下降。

针对问题五，从调峰压力、备用容量、电能质量等角度分析了高渗透率绿电园区对电力系统的影响，并提出了完善市场机制、配置储能、智能调度等政策建议。

本文的研究为绿电直连电氢氨园区的规划设计和优化运行提供了系统的数学建模方法和定量分析依据。
\end{abstract}
